\chapter*{Contribuciones Personales}
\label{cap:contribucionesPersonales}
\addcontentsline{toc}{chapter}{Contribuciones Personales}

\section*{David Cendejas Rodríguez}
Mis principales contribuciones a este trabajo se han centrado en el liderazgo del desarrollo técnico, la investigación de arquitecturas de Deep learning y la implementación de los diversos prototipos que han dado forma al sistema final, el prototipo 5. Mi labor ha abarcado desde el planteamiento arquitectónico inicial hasta el entrenamiento y validación de los modelos, asumiendo la responsabilidad de investigar las tecnologías de Machine Learning y Deep Learning aplicadas al ámbito del audio. Este proceso ha supuesto un reto constante, especialmente al abordar la síntesis FM, debido a problemas intrínsecos como la nula inyectividad y la extrema sensibilidad de sus parámetros, factores que han condicionado directamente la convergencia y el rendimiento de nuestro modelo de Sound Matching. 

Debido a ello, también he sido el principal responsable del desarrollo del contenido del Estado de la Cuestión de la memoria, así como realizando los diagramas necesarios para las partes más técnicas.

Paralelamente, he llevado a cabo una exploración técnica profunda sobre las representaciones de los datos y las estrategias de normalización. Esta parte del trabajo ha sido fundamental para transformar señales de audio en estructuras de datos procesables, como  tensores espectrograma, asegurando la estabilidad del entrenamiento. En este sentido, diseñé e implementé pipelines robustos para el preprocesamiento de audio, garantizando que la entrada al modelo fuera óptima para el aprendizaje de las características tímbricas y unificarlo a los procesos de inferencia y evaluación, permitiendo que el sistema se comportara y escalara de forma coherente en todas sus fases del desarrollo.

En lo relativo a la ingeniería de software, he sido el principal responsable de la implementación del código en Python, aplicando principios de Programación Orientada a Objetos para garantizar un sistema modular y escalable. Debido a la naturaleza investigadora y no lineal del proyecto, el código evolucionó de forma orgánica tras cada reunión con los tutores y la revisión de literatura técnica. Esto me exigió realizar constantes tareas de abstracción y refactorización para evitar la redundancia y mantener la encapsulación, asegurando siempre la eficiencia computacional y la legibilidad del software. Asimismo, gestioné íntegramente las fases de depuración y manejo de errores, lo que resultó crucial para asegurar la integridad de los resultados obtenidos en un entorno de desarrollo tan dinámico.

Finalmente, al disponer del hardware con mayor capacidad de cómputo del equipo, asumí la responsabilidad técnica de la fase de entrenamiento: Configuración de entornos de ejecución y gestión de recursos. Monitorización de métricas de evaluación estándar en la industria del Machine Learning para validar el rendimiento de los modelos en inferencia. Ajuste de hiperparámetros basado en la observación directa de las curvas de aprendizaje. Y entrenamiento de los modelos usados por mi compañero para realizar las pruebas y comparaciones.

\section*{Ángel Jiménez Izquierdo}
Al menos dos páginas con las contribuciones del estudiante 2. En caso de que haya más estudiantes, copia y pega una de estas secciones.

